% Full tail return propagation. Additive receiving proof; accepted providers retained.
\begingroup
\section{The full original tail through conductor transport and all receiving values}
\label{sec:FTR}
This calculation receives the explicit original relation and simultaneous
flag matrices NTG/SLC/STF/NFI into the original IPR, CPQ, NFR and CAI
objects. It preserves the programme's recorded origin and the distinct
creators and dates in its two deposited source notes
\cite{human:globalization2026,human:splitzero2026}.
The exact internal maps are stated with their source locators, while
the human mathematical dependencies are cited at their uses below.
All statements concern the original simple-quartet and fixed-period
domain specified in each section. No additional multiplicity or
asymptotic hypothesis is inserted.

\subsection{The original data and the two exact tail matrices}
Fix the actual simple quartet, its period and logarithm branch before
the degree limit. Retain
\[
 k=4l+1\ge9,\quad j=k+4,\quad q=(j+1)^2,\quad q'=(j-7)^2,
 \quad L=q-v,\quad m=L-q',\quad s\in\{1,j\}.                 \tag{FTR1}
\]
The complete conductor is
$E_A(z)=\sum_{u,w=0}^{8}a_{uw}e^{\beta_{8,u,w}z}$, with
$\beta_{a,u,w}=a/2+(2u-a)\delta+i(2w-a)\gamma$,
$0<\delta<1/2$, $\gamma>2$, and
$\mu_d=E_A^{(d)}(0)$, $v=\operatorname{ord}_0E_A$,
$\mu_v\ne0$. The relations, including their scalar factors, are
$g_r=\mathcal T_A(\chi_jS^r)$,
$\ell_r=\mu_v\binom{q+r}{v}$ and $u_r=g_r/\ell_r$.
Use the unchanged native line $S'=j/2-4+iy$ and measure
\[
 dm_s(y)=\frac{(2\pi)^{s/2}2^{s/2}}{4\pi\Gamma(s/2)}
             |\Gamma(s/4+iy/2)|^2\,dy .                  \tag{FTR2}
\]
The orthogonal polynomials, full norms and phases are exactly NTG3:
$p_{n,s}(y)=n!P_n^{(s/4)}(y/2;\pi/2)$ and
$h_{n,s}=(2\pi)^{s/2}n!(s/2)_n$. This is the stated
Meixner--Pollaczek convention, with its factor $dy=2\,dx$ and
mass checked in HPS1--4 and NTG3
\cite[eqs.~(18.19.8)--(18.19.9), (18.22.8)]{human:dlmf18}.
NTG4--12 and SLC3--20 evaluate every entry of the original graph
$X_r=[x_0,\ldots,x_r]$ and invertible functional matrix $A$.
In particular the positive polynomial modification in that entry
calculation is the complete confluent determinant of Szeg\H{o},
including the original leading cofactor and both rows at every root
\cite[Theorem 2.5, pp.~29--30]{human:szego1975}; the full substitution
and its polynomial proof are SLC1--21. None of those factors is changed
in this receiving calculation.

Let $\Omega_r=I_m+X_rX_r^*$, $\Omega_{-1}=I_m$, and keep any
integer $0\le R<q$. For every $R\le t\le q$ put
\[
\begin{gathered}
 G=\overline{\Omega_R},\quad H_1=A^*GA=D_{q+R},\quad
 H=\operatorname{diag}(H_1,4),\quad
 Y_t=[x_{R+1},\ldots,x_t],\quad Z_t=G^{-1/2}\overline{Y_t},\\
 \mathscr Z_t=\operatorname{diag}(Z_t,4),\qquad
 \mathscr N_t=\mathscr Z_t\mathscr Z_t^*,\qquad
 \mathcal U=\mathbb C^{4m} .                              \tag{FTR3}
\end{gathered}
\]
Here $D_N=D_{j,N-v}^{(s)}$ is the original native analytic quotient
form, and $\mathcal D_N=\operatorname{diag}(D_N,4)$.
For each fixed original flag $X\subset\mathcal U$ with independent
frame $X_0$, define
\[
\begin{gathered}
 \widehat Z_X=\operatorname{diag}(G^{1/2}A,4)
                 X_0(X_0^*HX_0)^{-1/2},\qquad
 \widehat P_X=\widehat Z_X\widehat Z_X^*,\\
 \widehat K_{X,t}=\mathscr N_t^{1/2}\widehat P_X
                                      \mathscr N_t^{1/2},\qquad
 f_t(X)=\log\det(I_{4m}+\widehat K_{X,t}).                 \tag{FTR4}
\end{gathered}
\]
An empty frame has projector zero and determinant one. The square
roots are the unique positive Hermitian square roots; all inverses
are on the indicated positive forms. The $\widehat P_X$ notation
in FTR4 means the explicit unitary change of flag coordinates in
NFI7, and does not denote the later conductor-transported metric.
The latter will have a superscript ${\rm tr}$ below.

For clarity the exact coordinate map behind FTR4 is
$U_4=\operatorname{diag}(G^{1/2}AH_1^{-1/2},4)$, with inverse
$U_4^*$. The row permutation and phases are NFI3's
$\mathscr R_t=\operatorname{diag}(\Theta_t^*,4)\Pi_t$,
$\Theta_t=\operatorname{diag}_{r=R+1}^{t}(-i^{-(L+r)})$.
NFI4 proves $\mathscr R_tC_t=\mathscr Z_t^*U_4$ directly from
$\rho_{L+r}/\sqrt{1+\zeta_{L+r}}=-i^{-(L+r)}x_r^{\mathsf T}A$.
Thus no phase or transpose is inferred from a determinant equality.
The range isometry in NFI9 identifies every nonzero eigenvalue in
FTR4, with multiplicity, with its original stacked-row compression;
its inverse is the adjoint on the image, and both complementary
kernels have zero contribution to the two logarithmic functions.
It follows in the same original frame that
\[
 f_t(X)=\log\frac{\det(X_0^*\mathcal D_{q+t}X_0)}
                              {\det(X_0^*HX_0)} .        \tag{FTR5}
\]
This is precisely STF6 and NFI8--9, now with all input entries
supplied by SLC16--21.

\subsection{Every original quotient and both allocations}
For a fixed original pair $K\subset Y$ choose its original kernel
frame $X_K$ and fixed lifts $X_Q$ of its quotient values. Set
\[
\begin{gathered}
 Z=X_Q-X_K(X_K^*HX_K)^{-1}X_K^*HX_Q,\qquad Q_0=Z^*HZ,\\
 E=H^{1/2}ZQ_0^{-1/2},\quad \widehat E=U_4E,\\
 Q_t=Q_0^{1/2}\left[I+\widehat E^*\mathscr Z_t
 (I+\mathscr Z_t^*\widehat P_K\mathscr Z_t)^{-1}
                  \mathscr Z_t^*\widehat E\right]Q_0^{1/2},\\
 \tau_{Y/K}=\sum_{t=q-1,q}\log\frac{\det Q_t}{\det Q_0}
            =\sum_{t=q-1,q}\{f_t(Y)-f_t(K)\}.             \tag{FTR6}
\end{gathered}
\]
For $K=0$ omit its kernel term; for $Y=K$ all quotient operators
are the unique operators on the zero space and its determinant is
one. These are the actual unchanged quotient values. Indeed every
representative of $x$ has transformed coordinate
$EQ_0^{1/2}x+Z_Ka$, where $Z_K$ frames $H^{1/2}K$
orthonormally. Its new norm is
$\|Q_0^{1/2}x\|^2+\|a\|^2+
\|C_tEQ_0^{1/2}x+C_tZ_Ka\|^2$.
Completing the square gives
$a=-(I+Z_K^*C_t^*C_tZ_K)^{-1}Z_K^*C_t^*C_tEQ_0^{1/2}x$.
Substitution followed by the actual unitary row map proves the
matrix in FTR6. In the fixed frame $[X_K,X_Q]$, block elimination
gives $\det\operatorname{Gram}(Y)=
\det\operatorname{Gram}(K)\det Q$, proving its determinant ratio.
This is the complex Hermitian form of the full Schur minimum and
factorization in Boyd and Vandenberghe, Appendix A.5.5, p.~650,
and Appendix C.4.1, p.~673 \cite{human:boyd-vandenberghe}; all
positive blocks, value fibres and cross terms were just specified.
In particular $Q_t\succeq Q_0$ and $\tau_{Y/K}\ge0$.

Use the original CPQ maps without choosing new subspaces:
\[
\begin{gathered}
 B=\operatorname{im}B_{\rm cof},\quad
 I=\operatorname{im}(B_{\rm cof}J_0),\quad
 S=\operatorname{im}[B_{\rm cof},F_e],\quad
 J=\ker\operatorname{diag}(B_j,4),\\
 K_i=I\cap J,\quad W_i=I+J,\quad W=B+J,\quad
 K_c=S\cap W,\quad T=S+J,\qquad e=0,1 .                 \tag{FTR7}
\end{gathered}
\]
The exact pairs and resulting tail returns are
\[
\begin{array}{c|c|c}
 \text{original form}&(Y,K)&\text{tail}\hline
 Q_b^E&(W,W_i)&\tau_b\!:=\tau_{W/W_i}\rule{0pt}{2.2ex}\\
 Q_{{\rm rem},e}&(\mathcal U,T)&\tau_{\rm rem}\!:=\tau_{\mathcal U/T}\\
 G_{s,e}&(S,K_c)&\tau_s\!:=\tau_{S/K_c}\\
 A_{\rm inv}&(I,K_i)&\tau_A\!:=\tau_{I/K_i}\\
 H_{\rm inv}&(W_i,J)&\tau_I\!:=\tau_{W_i/J}\\
 G_{t,e}&(T,W)&\tau_t\!:=\tau_{T/W}.
\end{array}                                                    \tag{FTR8}
\]
The first four are the same-value identifications CPQ7--10 and
STF14; the fifth is proved below on the fixed original values
$U_R$. The last is CPQ7's connecting target quotient. For each
pair the complete metric, including its original frame factors,
is FTR6. Only same-frame ratios annihilate these constant factors.

Define
\[
\begin{gathered}
 \tau_K={1\over4}\tau_{\mathcal U/J},\qquad
 \tau_R={1\over4}\sum_{t=q-1,q}f_t(J),\qquad
 \tau_{\rm total}=\log
       {\det\Omega_{q-1}\det\Omega_q\over(\det\Omega_R)^2},\\
 (X_1,\ldots,X_8)=(T,K_c,W_i,K_i,W,S,J,I),\quad
 (\sigma_1,\ldots,\sigma_8)=(1,1,1,1,-1,-1,-1,-1),\qquad
 \tau_L=\sum_{t=q-1,q}\sum_{a=1}^8\sigma_af_t(X_a).
                                                               \tag{FTR9}
\end{gathered}
\]
All eight positions are retained when subspaces coincide.
Substituting FTR8 and cancelling identical terms proves the
three simultaneous equalities
\[
\begin{gathered}
 \tau_L=\tau_t-\tau_s+\tau_I-\tau_A,\qquad
 4\tau_K=\tau_L+\tau_b+\tau_{\rm rem}+\tau_s+\tau_A,\\
 \tau_K+\tau_R=\tau_{\rm total}.                         \tag{FTR10}
\end{gathered}
\]
The last equality also follows from NFI6 applied at both endpoints:
$f_t(\mathcal U)=4\log(\det\Omega_t/\det\Omega_R)$.
The two allocations here are the original value quotient and
original kernel restriction, not selected fractions of their sum.

For a log determinant $h(N)$ in any fixed original value frame set
\[
\begin{aligned}
 \mathcal Rh&=-h(q-1)-h(q)+h(2q-1)+h(2q),\\
 \mathcal R_{\le R}h&=2h(q+R)-h(q)-h(q-1),\\
 \mathcal R_{>R}h&=h(2q-1)+h(2q)-2h(q+R).
\end{aligned}                                                   \tag{FTR11}
\]
They sum exactly, and they are the original weighted increments
with weights one at $r=0,q$ and two at $0<r<q$.
For each form in FTR8 write $p_Y=\mathcal R_{\le R}\log\det Q_Y$;
write $p_K,p_R$ for the original single-copy allocations and
$p_L$ for the signed original prefix. Then the complete native
returns have the exact expressions
\[
\begin{gathered}
 V_Y=p_Y+\tau_Y,\quad S_K=p_K+\tau_K,\quad
 S_R=p_R+\tau_R,\quad \mathcal L_e=p_L+\tau_L,\\
 4p_K=p_L+p_b+p_{\rm rem}+p_s+p_A,\quad
 p_K+p_R=p_{\rm total}.                                 \tag{FTR12}
\end{gathered}
\]
IPR3 proves telescoping and IPR4--8 proves these original fibre
identities at each part of the window. Positivity of the same
increment gives, with the literal NIB/IPR budget
$\mathfrak E_{s,R}=\sum_{r=0}^R\vartheta_rE_{j,s,r}$,
\[
\begin{gathered}
 |p_L|\le4\mathfrak E_{s,R},\qquad
 0\le p_Y\le4\mathfrak E_{s,R}\quad
       (Y=b,{\rm rem},s,A,I,t),\\
 0\le p_K,p_R\le\mathfrak E_{s,R},\qquad
 0\le p_K+p_R\le\mathfrak E_{s,R}.                       \tag{FTR13}
\end{gathered}
\]
For the two added pairs $I,t$, this bound follows by applying
the same complete minimum to their inclusions in the four-copy
ambient space. Explicitly, extend the fixed pair frame to a full
frame $C=[X_K,X_Q,X_{\rm out}]$. Two full Schur eliminations give
$|\det C|^2\det F=
\det(F|_K)\det Q_{Y/K,F}\det Q_{\mathcal U/Y,F}$
for each positive metric $F$. Under one positive update each
same-value factor is nondecreasing: restriction preserves the
quadratic inequality and minimization over the identical affine
fibre preserves it as well. Their determinant ratios therefore
are at least one and their product is the full ambient ratio.
Each quotient ratio is thus at most that ambient ratio. The latter
is $(1+a_r)^4$ by the original four-copy update. Taking logs
and the original weights proves FTR13. This argument retains
the original pair and does not substitute a space of equal rank.

\subsection{Directed conductor transport at the actual evaluation sites}
Let a superscript ${\rm tr}$ mean the identical coefficient and value
maps endowed with NCT's conductor-transported metric. For any one
restriction or attained quotient $Y$ of actual rank $r_Y$, write
\[
 d_Y(N)=\log\det Q_Y^{\rm tr}(N)-\log\det Q_Y(N),\qquad
 -l_Y(N)\le d_Y(N)\le u_Y(N).                            \tag{FTR14}
\]
Here $l_Y(N)=\mathcal L^{[4]}_{j,N-v,s,r_Y}$ and
$u_Y(N)=\mathcal U^{[4]}_{j,N-v,s,r_Y}$ are precisely
NCT19--20 at any fixed admissible original SGT radius. Explicitly,
with $h_4=\min(r_Y,4H_\varrho)$, they are
\[
\begin{aligned}
 u_Y(N)&=4r_YL_{s,N-v,v}
                 +2r_Y\log^+F_{\varrho,N-v},\\
 l_Y(N)&=4r_YL_{s,N-v,v}+2h_4\log^+G_{0,N-v}
                 +2(r_Y-h_4)\log^+V_{\varrho,N-v}^{\rm cap}.
\end{aligned}                                                   \tag{FTR15}
\]
Every radius, principal part, exceptional rank, residue and
diagonal factor in these already evaluated NCT constants is kept.
They are zero at rank zero. The transport is the actual complex-linear
functional map $\Lambda_F\mapsto\Lambda_{E_AF}$, in its ordinary-
transpose matrix, together with its complete lower quotient, as in
NCT11--14. In particular it applies to the same quotient values in
FTR6, rather than to arbitrarily chosen Euclidean flag frames.

For $Z\in\{\mathrm{full},\le R,>R\}$ let
$\eta_Y^Z=\mathcal R_Zd_Y$, using $\mathcal R$ for ``full''.
The following directed costs preserve all literal endpoint signs:
\[
\begin{array}{c|l|l}
 Z&E_Y^{Z,-}&E_Y^{Z,+}\hline
 \mathrm{full}
 &u_Y(q-1)+u_Y(q)+l_Y(2q-1)+l_Y(2q)
 &l_Y(q-1)+l_Y(q)+u_Y(2q-1)+u_Y(2q)\rule{0pt}{2.5ex}\\
 \le R&2l_Y(q+R)+u_Y(q)+u_Y(q-1)
      &2u_Y(q+R)+l_Y(q)+l_Y(q-1)\\
 >R&l_Y(2q-1)+l_Y(2q)+2u_Y(q+R)
    &u_Y(2q-1)+u_Y(2q)+2l_Y(q+R).
\end{array}                                                     \tag{FTR16}
\]
Then $-E_Y^{Z,-}\le\eta_Y^Z\le E_Y^{Z,+}$ and
$\eta_Y^{\rm full}=\eta_Y^{\le R}+\eta_Y^{>R}$ exactly.
Proof: substitute FTR14 in each term of FTR11. A negative coefficient
exchanges the roles of its two endpoint costs; coefficient two
multiplies its actual cost by two. This proves every table entry,
without adding per-row errors or assuming independent extremizers.
At repeated sites one may combine their actual coefficients first;
the displayed uncombined nonnegative costs remain valid.

For the signed eight-flag expression, write
$\eta_{X,\mathrm{res}}^Z$ for the residual of the restriction
form on the subspace $X$, and put
$\eta_L^Z=\sum_a\sigma_a\eta_{X_a,\mathrm{res}}^Z$ and
\[
 E_L^{Z,-}=\sum_{\sigma_a=1}E_{X_a}^{Z,-}
              +\sum_{\sigma_a=-1}E_{X_a}^{Z,+},\quad
 E_L^{Z,+}=\sum_{\sigma_a=1}E_{X_a}^{Z,+}
              +\sum_{\sigma_a=-1}E_{X_a}^{Z,-}.          \tag{FTR17}
\]
The $E_{X_a}^{Z,\pm}$ in this display are the costs for those
restriction forms. These bounds keep all eight labelled occurrences.
Define
$\eta_{\rm tar}^Z=4(S_K^{Z,\rm tr}-S_K^Z)$ by applying FTR14--16
directly to the full four-copy target quotient of rank $4t_j$;
define $\eta_J^Z=\eta_{J,\mathrm{res}}^Z
 =4(S_R^{Z,\rm tr}-S_R^Z)$ by applying them
directly to the actual four-copy kernel restriction. The transported
and native determinant factorizations at the same sites give
\[
\begin{gathered}
 \eta_{\rm tar}^Z=\eta_L^Z+\eta_b^Z+\eta_{\rm rem}^Z+\eta_s^Z+\eta_A^Z,
 \qquad \eta_{\rm tar}^Z+\eta_J^Z=\eta_{\mathcal U,\mathrm{res}}^Z,\\
 \eta_L^Z=\eta_t^Z-\eta_s^Z+\eta_I^Z-\eta_A^Z .         \tag{FTR18}
\end{gathered}
\]
Here $\eta_I$ is the quotient residual for $H_{\rm inv}$ in
FTR8; it is distinct from the restriction residual
$\eta_{I,\mathrm{res}}$. Similarly $\eta_{\rm tar}$ denotes
the full target quotient, while $\eta_{T,\mathrm{res}}$ is
one of the eight restriction residuals. These are equalities
of the actual residuals. For example their
first line follows by subtracting the full CPQ identity for the
native metric from that for the transported metric and applying
the same $\mathcal R_Z$. Thus the direct target interval in
FTR16 can be intersected with the interval supplied by its first
identity; their extremizers have not been chosen separately.

The resulting complete returns are explicitly
\[
\begin{gathered}
 V_Y^{\rm tr}=\tau_Y+p_Y+\eta_Y^{\rm full},\qquad
 \mathcal L_e^{\rm tr}=\tau_L+p_L+\eta_L^{\rm full},\\
 4S_K^{\rm tr}=4\tau_K+4p_K+\eta_{\rm tar}^{\rm full},\qquad
 4S_R^{\rm tr}=4\tau_R+4p_R+\eta_J^{\rm full} .          \tag{FTR19}
\end{gathered}
\]
In particular each $V_Y^{\rm tr}$ lies in the finite interval
$[\tau_Y-E_Y^{\mathrm{full},-},
\tau_Y+4\mathfrak E_{s,R}+E_Y^{\mathrm{full},+}]$,
intersected with $[0,\infty)$. The direct allocation intervals
are obtained by using $4\mathfrak E_{s,R}$ and dividing their
four-copy endpoints by four. The exact formulas FTR18--19 retain
the correlation stronger than these separate enclosures.

\subsection{The complete transported threshold profile from the common matrices}
Define the original finite signed count
\[
 \mathcal B_{e,s}(R)=\sum_{t=q-1,q}\int_1^\infty
       \sum_{a=1}^8\sigma_a
       \#\{\lambda_b(\widehat K_{X_a,t})>u\}\,{du\over u}.
                                                               \tag{FTR20}
\]
Every matrix is the full original-data FTR3--4 matrix. Each integral
terminates past its finite largest eigenvalue. Set $h(x)=
\log(1+x)-\log\max(1,x)$ and retain the actual finite residual
\[
 \epsilon_{e,s}(R)=\sum_{t=q-1,q}\sum_{a=1}^8\sigma_a
                 \sum_b h(\lambda_b(\widehat K_{X_a,t})).
                                                               \tag{FTR21}
\]
The elementary identity
$\log\max(1,\lambda)=\int_1^\infty\mathbf1_{u<\lambda}\,du/u$
proves $\tau_L=\mathcal B_{e,s}(R)+\epsilon_{e,s}(R)$ exactly.
Let $d_t=\operatorname{rank}\mathscr N_t$. The two actual nested
chains $J\subset W_i\subset W\subset T$ and
$K_i\subset I\subset K_c\subset S$ imply
\[
 |\epsilon_{e,s}(R)|\le
       \varepsilon_R:=2(d_{q-1}+d_q)\log2\le16m\log2.    \tag{FTR22}
\]
Here is the scalar error argument. Write $h=u-v$ with
$u(x)=\log(1+\min(x,1))$ and $v(x)=0$ for $x\le1$,
$v(x)=\log(2x/(1+x))$ for $x\ge1$. Both are increasing
from zero and bounded by $\log2$. For four ordered eigenvalues,
the sum of the first and last disjoint increments of each function
is between zero and $\log2$, by adding its middle nonnegative
increment. Their difference lies between $-\log2$ and $\log2$.
The two flag chains therefore cost $2\log2$ for each of the common
$d_t$ range coordinates. Matrix order gives the eigenvalue order
by the minimum--maximum characterization of Hermitian eigenvalues,
whose elementary proof is included in NCT18. This proves FTR22,
including singular matrices and coincident flags.

Let $\mathcal P_e^{\rm tr}$ be the original transported successive-row
threshold sum NPK10, with its original endpoint weights. Define its
actual difference from its own determinant return by
$\chi_e^{\rm tr}=\mathcal P_e^{\rm tr}-\mathcal L_e^{\rm tr}$.
At each row the same argument uses four eigenvalues and two chains,
so its cost is $8\log2$. The original weights sum to $2q$, hence
\[
 |\chi_e^{\rm tr}|\le16q\log2,\qquad
 \boxed{\mathcal P_e^{\rm tr}
 =\mathcal B_{e,s}(R)+p_L+\epsilon_{e,s}(R)
                       +\eta_L^{\rm full}+\chi_e^{\rm tr}.}    \tag{FTR23}
\]
This compares the complete transported row sum directly with the
native simultaneous matrices. In particular its directed enclosure is
\[
\begin{split}
 -4\mathfrak E_{s,R}-\varepsilon_R-E_L^{\mathrm{full},-}
                       -16q\log2
 &\le \mathcal P_e^{\rm tr}-\mathcal B_{e,s}(R)\\
 &\le4\mathfrak E_{s,R}+\varepsilon_R+E_L^{\mathrm{full},+}
                       +16q\log2 .                      \tag{FTR24}
\end{split}
\]
The bound for $\mathcal L_e^{\rm tr}-\mathcal B_{e,s}(R)$
omits the last $16q\log2$ on each side. The original native
threshold-to-determinant error is not paid a second time in FTR24.
Each residual in FTR23 has its specified actual expression;
there is no assertion that individual endpoint and row eigenvalues
are equal. Replacing ``full'' by either partial operator in
FTR18--19 gives the corresponding complete partial identities.

At the already established $R_j=\lfloor j^{3/2}/\log j\rfloor$
with $R_j<q$, IPR9 gives
$\mathfrak E_{s,R_j}=O_{\rm actual}(j^{5/2}+j^3/\log j)
=o(kq_k)$. The finite transport bounds may be replaced by the
infimum of their valid endpoint intervals over admissible radii;
equivalently use NPK11's absolute bounds for the asymptotic statement.
For each fixed radius NCT21 first takes the degree limit, and only
afterwards admissible radii tend to one. Since all actual flag and
quotient ranks here are $O(j)$, this yields $o(kq_k)$ for every
four-site or three-site radius-infimum transport cost. No fixed-radius
cost is asserted to be little-oh. Since $m=O(j)$ and $q=O(j^2)$,
FTR22--24 prove for each original order and primitive index that
\[
 \mathcal P_e^{\rm tr}-\mathcal B_{e,s}(R_j)=o(kq_k),\qquad
 \mathcal L_e^{\rm tr}-\mathcal B_{e,s}(R_j)=o(kq_k).       \tag{FTR25}
\]
The exact quantity on the right of FTR20 remains the original
correlated spectral count. This proof evaluates the full receiving
map and finite costs; it assigns no new asymptotic coefficient to
that count. The independently retained proper source and complete
arithmetic action are treated next through their actual value maps.

All prefixes in this receiver also have an explicit finite graph
formula, rather than merely the bounds FTR13. Set
$\mathcal A=\operatorname{diag}(A,4)$ and
$\mathcal O_r=\operatorname{diag}(\overline{\Omega_r},4)$.
For every unchanged restriction frame $X_0$, put
\[
 p_{X,\mathrm{res}}=2\log\det(X_0^*\mathcal A^*\mathcal O_R\mathcal A X_0)
       -\log\det(X_0^*\mathcal A^*\mathcal O_0\mathcal A X_0)
       -\log\det(X_0^*\mathcal A^*\mathcal A X_0).        \tag{FTR26}
\]
The target cutoffs $q+R,q,q-1$ correspond to native graph
indices $R,0,-1$, respectively, so NTG20 and FTR11 prove this
formula including $R=0$. Every quotient prefix is the difference
$p_{Y,\mathrm{res}}-p_{K,\mathrm{res}}$ in its fixed pair;
in particular $p_I$ in FTR12--13 equals
$p_{W_i,\mathrm{res}}-p_{J,\mathrm{res}}$.
Every signed prefix is the original
eight-term combination. Frame constants cancel because the
coefficients in FTR26 sum to zero. Thus every prefix in FTR19
and FTR23 is itself evaluated in the same original finite coefficients.

For $R=R_j$ and $t=q-1,q$ let $\mathfrak U_{R,t}$ be the
minimum of the three explicitly evaluated bounds NTG31:
\[
\begin{split}
 \mathfrak U_{R,t}=\min\bigg\{&\sum_{a=R+1}^{t}E_{j,s,a},\quad
 d_0\log\left(1+{\sum_{a=R+1}^{t}x_a^*\Omega_R^{-1}x_a\over d_0}\right),\\
 &\sum_{a=R+1}^{t}\log(1+x_a^*\Omega_R^{-1}x_a)
       -\Delta_{R,t}\bigg\},\qquad d_0=\min(m,t-R),
                                                               \tag{FTR27}
\end{split}
\]
where $\Delta_{R,t}$ is the full adjacent-pair correction NTG28,
including each actual $x_a^*\Omega_R^{-1}x_b$. At an empty endpoint
set $\mathfrak U_{R,R}=0$. Define
$\mathfrak U_R=\mathfrak U_{R,q-1}+\mathfrak U_{R,q}$.
The complete Schur determinant NTG25 proves
$0\le\tau_{\rm total}\le\mathfrak U_R$.
Each quotient ratio is at most its four-copy ambient ratio,
as proved for FTR13, and each allocation is nonnegative with
sum $\tau_{\rm total}$. The original two-chain signed estimate
NIB16 therefore yields simultaneously
\[
\begin{gathered}
 0\le\tau_K,\tau_R\le\mathfrak U_R,\qquad
 0\le\tau_Y\le4\mathfrak U_R\quad
                 (Y=b,{\rm rem},s,A,I,t),\\
 |\tau_L|\le4\mathfrak U_R,\qquad
 |\mathcal B_{e,s}(R)|\le4\mathfrak U_R+\varepsilon_R .  \tag{FTR28}
\end{gathered}
\]
For the signed bound, the same-chain disjoint-increment argument
applied to $\log(1+x)$ at an original row bounds the absolute
signed increment by $4\log(1+a_r)$; summing the tail and using
NTG32 proves it directly. The count bound follows from FTR21--22.
These estimates may be inserted in the directed intervals FTR19
and FTR24 without deleting any of their correlated residuals.
No assertion that $\mathfrak U_R=o(kq_k)$ is used here.

Finally retain both original source orders, with a subscript $1$
or $j$ on each actual matrix and residual. For every positive
quotient in FTR8 the exact order difference is
\[
 V_{Y,j}^{\rm tr}-V_{Y,1}^{\rm tr}
 =(\tau_{Y,j}-\tau_{Y,1})+(p_{Y,j}-p_{Y,1})
                  +(\eta_{Y,j}^{\rm full}-\eta_{Y,1}^{\rm full}).
                                                               \tag{FTR29}
\]
Its excess above the displayed tail difference lies in
\[
\begin{split}
 [&-4\mathfrak E_{1,R}-E_{Y,j}^{\mathrm{full},-}
                             -E_{Y,1}^{\mathrm{full},+},\\
  &\phantom{-}4\mathfrak E_{j,R}+E_{Y,j}^{\mathrm{full},+}
                             +E_{Y,1}^{\mathrm{full},-}].        \tag{FTR30}
\end{split}
\]
This follows by subtracting the two actual FTR19 equalities,
using nonnegative prefixes FTR13 and directed errors FTR16.
For the signed profile use its two actual FTR23 identities instead;
the prefix difference has absolute bound
$4(\mathfrak E_{j,R}+\mathfrak E_{1,R})$, its two simultaneous
errors cost $\varepsilon_{R,j}+\varepsilon_{R,1}$, its directed
transport costs have the same subtracted signs as FTR30, and
the two original threshold residuals cost $32q\log2$.
This propagates the new matrices through the earlier two-order
receivers without treating a known difference as an absolute value.

\endgroup

\begingroup
\section{The original invariant value fibre and its full proper-source return}
\label{sec:PRD}

This calculation propagates the original ITL1--6, NFR1--12,
CPQ2--20 and EFP19,28--31 through the simultaneous matrices
STF4--18 and the exact relation-graph map NFI1--10. The classical
quadratic minimum and block determinant identity used at each quotient
are those in Boyd and Vandenberghe, Appendix A.5.5, p.~650,
equations (A.13)--(A.14), and Appendix C.4.1, p.~673,
equations (C.4)--(C.5) \cite{human:boyd-vandenberghe}.
The complex Hermitian calculation, every fibre and every map needed
here are given below. The original programme maps and results remain
attributed to their cited ITL/NFR/CPQ/EFP and NTG/STF/NFI providers;
the classical citation supplies the finite linear algebra operation.

\subsection{The unchanged domain and actual quotient values}
Retain the original hypothetical simple quartet, five-orbit period,
logarithm branch, complete nonreduced conductor $E_A$, its order
$v=\operatorname{ord}_0E_A$ and the full nonzero derivative
$\mu_v=E_A^{(v)}(0)$. The data and original rank-valid domain are
those of ITL1--3 and CPQ1--3, including the actual period-domain
condition $|u|\ge R_*$ wherever it is required by those maps. Put
\[
\begin{gathered}
 k=4l+1\ge9,\qquad j=k+4,\qquad
 q_a=(a+1)^2,\quad q'_a=(a-7)^2,\quad m_a=q_a-v-q'_a,\\
 c'_a=a/2-4,\quad \alpha_A=v!/\mu_v,\qquad
 p=0:(s_0,s_1)=(1,1),\quad p=1:(s_0,s_1)=(k,j),\\
 q=q_j,\qquad N\in[q-1,2q]\cap\mathbb Z,\qquad D=N-v.
\end{gathered}                                                    \tag{PRD1}
\]
The simple-quartet domain fixes $q_a$ as displayed; no change of
quartet multiplicity is made. Its conductor order $v$ is retained
even when the selected period makes it nonzero. Every source-order
identity below holds separately for $p=0,1$. Each connecting identity
holds separately for both primitive indices $e=0,1$.

Use the full original coefficient multiplier, maps and spaces
\[
\begin{gathered}
 \mathcal V=\mathbb C^{m_k},\quad
 \mathcal U=\mathbb C^{4m_j},\quad
 V=\mathbb C^{j_{0,k}},\qquad J_0:V\longrightarrow\mathcal V,\\
 \mathsf T_a=\alpha_A N_aZ_{\exp,a},\quad
 M_i=\mathsf T_j^{-1}m_{i,k}\mathsf T_k,\quad
 m_{i,k}[f]=[x_i^4f],\quad
 B_{\rm cof}=\operatorname{stack}_{i=1}^4M_i,\\
 F=B_{\rm cof}J_0:V\longrightarrow\mathcal U,\qquad
 \mathbb B=\operatorname{diag}(B_j,4):\mathcal U
          \longrightarrow\mathbb C^{4t_j},\qquad
 J=\ker\mathbb B,\\
 X=\mathbb BF,\quad E=\ker X,\quad
 I=\operatorname{im}F,\quad K_i=I\cap J,\quad W_i=I+J.
\end{gathered}                                                    \tag{PRD2}
\]
The maps $F$ and $J_0$ are injective and $B_j$ is onto by the
original ITL2--3/CPQ2 proof. Their ordinary-transpose coefficient
duality is unchanged; stars below are complex adjoints only in
Hermitian metric expressions. In particular the original maps
$N_a,Z_{\exp,a}$ and the multiplier $m_{i,k}$ remain in PRD2.

Choose exactly the original fixed ITL4 frame $Z=[Z_E,Z_R]$ of $V$,
where $Z_E$ spans $E$. Write
\[
 e_E=\dim E,\quad r_I=\dim V-e_E,
 \qquad F_E=FZ_E,\quad F_R=FZ_R,\quad
 U_R=XZ_R=\mathbb BF_R.                                  \tag{PRD3}
\]
The frame $F_R$ is the original coefficient lift of $U_R$.
The frame $U_R$ is injective:
if $XZ_Rx=0$, then $Z_Rx\in E$, while the complement property
of $Z$ forces $x=0$. Also
\[
 F E=K_i,\qquad
 I=F E\mathbin{\oplus}\operatorname{im}F_R,
 \qquad W_i=J\mathbin{\oplus}\operatorname{im}F_R.       \tag{PRD4}
\]
Indeed $Fa\in J$ is equivalent to $Xa=0$, which is $a\in E$.
This proves the first equality. Decomposing $a=Z_Eb+Z_Rx$
proves the second, using injectivity of $F$; adding $J$ proves
the third, and its directness follows by applying $\mathbb B$.

For complete value maps let $L_R=(U_R^*U_R)^{-1}U_R^*$,
the fixed left inverse on $\operatorname{im}U_R$. The two maps
and their inverses are
\[
\begin{aligned}
 I/K_i&\longrightarrow\mathbb C^{r_I},&
 [i]&\longmapsto L_R\mathbb Bi,&
 x&\longmapsto[F_Rx]\quad\text{(inverse)},\\
 W_i/J&\longrightarrow\mathbb C^{r_I},&
 [w]&\longmapsto L_R\mathbb Bw,&
 x&\longmapsto[F_Rx]\quad\text{(inverse)}.
\end{aligned}                                                     \tag{PRD5}
\]
These rules are well defined because $\mathbb B$
kills the indicated kernels; PRD4 proves that they are mutual
inverses. Composing either with $x\mapsto U_Rx$ identifies its
values with the very same original invariant image in
$\mathbb C^{4t_j}$. The inclusion $I\hookrightarrow W_i$
therefore induces the exact isomorphism
\[
 \iota_I:I/K_i\longrightarrow W_i/J,\qquad [i]\longmapsto[i].
                                                               \tag{PRD6}
\]
Its inverse sends $[i+j]$ to $[i]$: if two decompositions represent
the same class, their $I$ components differ by an element of
$I\cap J=K_i$. Both maps are the identity on the specified
$\mathbb C^{r_I}$ values. This proves their relationship before
any metrics or determinant ratios are introduced.

\subsection{All three full metrics and their attained lifts}
Put
\[
 \mathcal D_N^p=\operatorname{diag}(D_{j,N-v}^{(s_1)},4),
 \qquad G_N^p=J_0^*D_{k,N-v}^{(s_0)}J_0.                  \tag{PRD7}
\]
Each $D_{a,D}^{(s)}$ is the original positive coefficient metric
after its entire lower-evaluation minimum. More explicitly, with
the original lower and analytic-numerator rows through degree $D$,
\[
 D_{a,D}^{(s)}=T_{a,D}^*
 \bigl[I-V_{a,D}(V_{a,D}^*V_{a,D})^{-1}V_{a,D}^*\bigr]T_{a,D}.
                                                               \tag{PRD8}
\]
This is the MNT/DNE form used in ITL1 and STF2. Its rows retain
the original centre $c'_a$, conductor division moments and phases,
and the full squared norms
$h_{n,s}=(2\pi)^{s/2}n!(s/2)_n$. PRD8 requires the original
degree-$k$ rows when $a=k$ and degree-$j$ rows when $a=j$.
Positivity and full lower rank on PRD1 are the original MNT/DNE
rank results; no new rank assumption is introduced.

Here is the precise complex Hermitian version of the classical
minimum used below \cite[Appendix A.5.5, p.~650]{human:boyd-vandenberghe}.
For a positive Hermitian $P$, an independent kernel frame $K_0$
and value lifts $L_0$ independent modulo that kernel, set
\[
 \begin{split}
 \operatorname{Sch}(P;K_0,L_0)
  &=L_0^*PL_0-L_0^*PK_0(K_0^*PK_0)^{-1}K_0^*PL_0,\\
 z_P(x)&=L_0x-K_0(K_0^*PK_0)^{-1}K_0^*PL_0x.
 \end{split}                                                     \tag{PRD9}
\]
Every representative is $L_0x+K_0a$. Its squared norm is
\[
 (a+(K_0^*PK_0)^{-1}K_0^*PL_0x)^*(K_0^*PK_0)
 (a+(K_0^*PK_0)^{-1}K_0^*PL_0x)
 +x^*\operatorname{Sch}(P;K_0,L_0)x.
\]
The first term has a unique zero, proving PRD9's attained minimum
and lift. A triangular change of determinant one from
$[K_0,L_0]$ to $[K_0,z_P]$ diagonalizes its Gram by blocks;
therefore its determinant is
$\det(K_0^*PK_0)\det\operatorname{Sch}(P;K_0,L_0)$.
This proves the complex form of the cited Schur operation with
the actual affine fibre and all cross terms retained.

Choose any fixed independent frame $J_*$ for the entire $J$.
The three forms on the same original $Z_R$ values are exactly
\[
\begin{aligned}
 Q_N^p&=\operatorname{Sch}(G_N^p;Z_E,Z_R),\\
 A_N^p&=\operatorname{Sch}(\mathcal D_N^p;F_E,F_R),\\
 H_N^p&=\operatorname{Sch}(\mathcal D_N^p;J_*,F_R)
      =U_R^*(\mathbb B(\mathcal D_N^p)^{-1}\mathbb B^*)^{-1}U_R.
\end{aligned}                                                    \tag{PRD10}
\]
Thus $A_N^p$ is precisely $A_{\rm inv}^p$ for $(I,K_i)$,
and $H_N^p$ is precisely $H_{\rm inv}^p$ for $(W_i,J)$.
To prove the last equality directly, for a target value $y$ the
vector
\[
 z=(\mathcal D_N^p)^{-1}\mathbb B^*
       (\mathbb B(\mathcal D_N^p)^{-1}\mathbb B^*)^{-1}y
\]
satisfies $\mathbb Bz=y$ and is $\mathcal D_N^p$-orthogonal
to $J$. Every other lift differs from it by $J$, so its norm
is the full minimum. Taking $y=U_Rx$ proves PRD10 and agrees
with the third attained lift in PRD9. Because $K_i\subset J$,
the two target-order minima also give $H_N^p\preceq A_N^p$.
The same-valued metric map in PRD6 is exactly this contraction;
the independent source metric is $Q_N^p$ with the first
formula in PRD10.

In particular, every entry of $Q_N^p$ is the finite original
degree-$k$ expression
\[
\begin{split}
 Q_N^p={}&Z_R^*J_0^*D_{k,N-v}^{(s_0)}J_0Z_R\\
 &-Z_R^*J_0^*D_{k,N-v}^{(s_0)}J_0Z_E
 (Z_E^*J_0^*D_{k,N-v}^{(s_0)}J_0Z_E)^{-1}
 Z_E^*J_0^*D_{k,N-v}^{(s_0)}J_0Z_R .
\end{split}                                                      \tag{PRD11}
\]
For $E=0$ every term containing its frame is absent. For $r_I=0$
all value forms are the unique operators on the zero space,
their determinants are $1$ and all value returns below are $0$.
These conventions apply throughout, including PRD5's empty maps.

\subsection{Literal two-endpoint matrices on those same values}
Fix an original integer $0\le R<q$, put $N_*=q+R$ and
$\mathcal H=\mathcal D_{N_*}^p$. Use the already calculated
NFI matrices
\[
 \mathscr Z_t=\operatorname{diag}
 (\overline{\Omega_R}^{-1/2}\overline{Y_t},4),\quad
 Y_t=[x_{R+1},\ldots,x_t],\qquad t=q-1,q,                 \tag{PRD12}
\]
including the empty tail at $R=t=q-1$. These $x_r,\Omega_R$
are the full NTG relation coefficients at target degree $j$ and
source order $s_1$; the exact phase/permutation maps NFI3--4 are
retained. Let $U_4$ be the NFI1 unitary. For either original
pair $(Y,K)=(W_i,J)$ or $(I,K_i)$, use the value lifts $F_R$
and respectively the kernel frames $J_*$ or $F_E$, and define
\[
\begin{gathered}
 Z_{Y/K}=F_R-K_*(K_*^*\mathcal H K_*)^{-1}K_*^*\mathcal H F_R,
 \qquad Q_{Y/K,*}=Z_{Y/K}^*\mathcal H Z_{Y/K},\\
 E_{Y/K}=\mathcal H^{1/2}Z_{Y/K}Q_{Y/K,*}^{-1/2},\qquad
 \widehat E_{Y/K}=U_4E_{Y/K},\qquad
 \widehat P_K=U_4P_KU_4^*,\\
 \mathcal J_{Y/K,t}=I_{r_I}+
 \widehat E_{Y/K}^*\mathscr Z_t
 (I+\mathscr Z_t^*\widehat P_K\mathscr Z_t)^{-1}
 \mathscr Z_t^*\widehat E_{Y/K} .
\end{gathered}                                                    \tag{PRD13}
\]
The original $P_K$ is STF5's projector onto $\mathcal H^{1/2}K$.
Every displayed kernel term is omitted when the kernel is zero.
The same complete quadratic minimum PRD9, now after stacking the
actual tail rows, gives STF10; conjugating its row space by the
explicit unitary NFI3 gives
\[
 H_{q+t}^p=(H_{N_*}^p)^{1/2}\mathcal J_{W_i/J,t}(H_{N_*}^p)^{1/2},
 \quad
 A_{q+t}^p=(A_{N_*}^p)^{1/2}\mathcal J_{I/K_i,t}(A_{N_*}^p)^{1/2}.
                                                               \tag{PRD14}
\]
For clarity, the minimization giving the middle factor has the
literal form $\min_a\{\|a\|^2+\|y+Ba\|^2\}$, with
$B=C_t Z_K$ and $y=C_tE_{Y/K}Q_{Y/K,*}^{1/2}x$.
Here $C_t$ is the STF4 stack of the actual original row updates,
right-multiplied by $\mathcal H^{-1/2}$, and $Z_K$ is any
orthonormal column frame of $\mathcal H^{1/2}K$.
Its minimizer is $-(I+B^*B)^{-1}B^*y$; multiplication proves
$I-B(I+B^*B)^{-1}B^*=(I+BB^*)^{-1}$, yielding PRD13.
The NFI3 row conjugation changes $C_t$ into $\mathscr Z_t^*U_4$.
This proves PRD14 with every cross term, including singular tail
matrices. Its $Q_{Y/K,*}^{1/2}$ factors remain on both sides in
the original quotient-value space; they are not scalar masses.

Let $F_t(Y)$ retain exactly STF5/NFI8's restriction log ratio.
Applying the block determinant identity proved after PRD9 to
the fixed frame $[K_*,F_R]$ at the two cutoffs gives
\[
\begin{aligned}
 \log\det\mathcal J_{W_i/J,t}&=F_t(W_i)-F_t(J),\\
 \log\det\mathcal J_{I/K_i,t}&=F_t(I)-F_t(K_i).
\end{aligned}                                                    \tag{PRD15}
\]
Every fixed frame determinant cancels in its own two-cutoff
ratio. Define the two exact tail values
\[
 \tau_I^p=\sum_{t=q-1,q}\log\det\mathcal J_{W_i/J,t},\qquad
 \tau_A^p=\sum_{t=q-1,q}\log\det\mathcal J_{I/K_i,t}.
                                                               \tag{PRD16}
\]
They are nonnegative because PRD13 is the identity plus a positive
semidefinite matrix. This is the same value for the target invariant
form whether obtained by its complete $J$ minimum or by its
restriction to the original image $U_R$.

\subsection{The full relative return, with the degree-k source intact}
For any value form $Z_N$ let
\[
 \Pi_Z^p=2\log\det Z_{N_*}^p-\log\det Z_q^p
                            -\log\det Z_{q-1}^p.
                                                               \tag{PRD17}
\]
Use this definition for $Z=H,A$, writing their prefixes as
$\Pi_I^p,\Pi_A^p$. Direct addition of PRD14's two logarithmic
determinants and PRD17 gives
\[
 V_I^p=\Pi_I^p+\tau_I^p,\qquad
 V_A^p=\Pi_A^p+\tau_A^p.                                \tag{PRD18}
\]
These are the original four-endpoint returns with signs
$(-1,-1,+1,+1)$ at $N=(q-1,q,2q-1,2q)$, not new windows.
Define the proper-source term directly from PRD11 at those four
original arguments:
\[
 V_Q^p=\log\frac{\det Q_{2q-1}^p\det Q_{2q}^p}
                         {\det Q_{q-1}^p\det Q_q^p},
 \qquad \mathscr R_{X,N}^p=(Q_N^p)^{-1/2}H_N^p(Q_N^p)^{-1/2}.
                                                               \tag{PRD19}
\]
Thus at each upper endpoint its complete relative positive operator is
\[
 \mathscr R_{X,q+t}^p
 =(Q_{q+t}^p)^{-1/2}(H_{N_*}^p)^{1/2}
 \mathcal J_{W_i/J,t}(H_{N_*}^p)^{1/2}(Q_{q+t}^p)^{-1/2}.
                                                               \tag{PRD20}
\]
The order of these generally noncommuting matrices is part of the
formula. Taking determinants, with the exact inverse square root
on both sides, proves the full finite receiving identities
\[
\begin{aligned}
 \mathcal M^p&=\mathcal R\log\det\mathscr R_X^p
                 =\Pi_I^p+\tau_I^p-V_Q^p,\\
 \mathcal R\log\frac{\det A_{\rm inv}^p}{\det Q_{\rm inv}^p}
               &=\Pi_A^p+\tau_A^p-V_Q^p,\\
 \mathcal R\log\omega_{\rm inv}^p
               &=(\Pi_I^p-\Pi_A^p)+(\tau_I^p-\tau_A^p).
\end{aligned}                                                    \tag{PRD21}
\]
The first line is the NFR matched return, the second its actual
same-value source correction, and the third the original last
target contraction. Their equality uses the proved maps PRD5--6,
not an identification of the degree-$k$ and degree-$j$ metrics.

The proper source in every term of PRD19 is evaluated at
\[
 (D_0,D_1,D_2,D_3)=(q_j-v-1,q_j-v,2q_j-v-1,2q_j-v).
                                                               \tag{PRD22}
\]
Its degree is $k$, centre $c'_k=k/2-4$, and order $s_0=1$ or $k$.
Writing $b=q_j-q_k=8k+24$ proves that its two low cutoffs are
the standard degree-$k$ low cutoffs plus $b$, and the two high
cutoffs are the standard degree-$k$ high cutoffs plus $2b$.
The target matrices in PRD12--16 instead retain degree $j$,
centre $c'_j=j/2-4$ and order $s_1=1$ or $j$. Consequently
PRD21 keeps the exact same $Z_R$ values while explicitly
retaining both original polynomial row families and all four
source displacements. No estimate for $V_Q^p$ has been inserted
in this identity.

The exact target-prefix cost can also be retained as a finite interval.
For any original $K\subset Y\subset\mathcal U$, restriction and
minimum over the fixed fibre preserve the increasing metric
$\mathcal D_N^p$. Complete a fixed frame for $K$, then the
specified lifts for $Y/K$, then lifts for $\mathcal U/Y$.
Successive Schur factorization expresses the ambient determinant
increment as the sum of three nonnegative log determinant
increments. Thus the quotient increment lies between zero and
the ambient increment. For either pair in PRD13 the latter is
$4\log(1+a_r)$ by the four-copy original IPR1 update. With
the original weights $\vartheta_0=\vartheta_q=1$ and
$\vartheta_r=2$ inside the window, this proves
\[
 0\le\Pi_I^p,\Pi_A^p
 \le4\sum_{r=0}^R\vartheta_r\log(1+a_r)
 \le4\mathfrak E_{s_1,R}.                              \tag{PRD23}
\]
The last finite constant is exactly the already proved IPR2/NIB14
budget. In particular, writing $Z=I$ or $A$ and interpreting its
left side respectively as $\mathcal M^p$ or
$\mathcal R\log(\det A_{\rm inv}^p/\det Q_{\rm inv}^p)$,
the full finite interval is
\[
 \tau_Z^p-V_Q^p\ \le\ \Pi_Z^p+\tau_Z^p-V_Q^p
       \ \le\ \tau_Z^p-V_Q^p+4\mathfrak E_{s_1,R}.       \tag{PRD24}
\]
It has the original independent proper-source term on both sides.
At the already established $R_j=\lfloor j^{3/2}/\log j\rfloor$
on its original $R_j<q$ domain, this width is $o(kq_k)$ by
IPR9. This bounds the removed target prefix; it does not assign
an asymptotic value to either retained term.

\subsection{The connecting receiver and all compatibility terms}
Use precisely the STF13/CPQ flags $S,K_c,W,T$ and the original
connecting columns at primitive index $e$. Their quotient pairs
for $G_t,G_s,A_{\rm inv},H_{\rm inv}$ are respectively
$(T,W),(S,K_c),(I,K_i),(W_i,J)$. Applying PRD15 to these pairs
and substituting the actual eight flags of STF16 gives, by
literal cancellation of all eight labelled occurrences,
\[
 \mathcal L_e^{p,>R}
  =\tau_I^p+\tau_t^p-\tau_s^p-\tau_A^p,\qquad
 \mathcal L_e^{p,\le R}
  =\Pi_I^p+\Pi_t^p-\Pi_s^p-\Pi_A^p.                    \tag{PRD25}
\]
Here $\tau_t=\sum_{h=q-1,q}[F_h(T)-F_h(W)]$ and
$\tau_s=\sum_{h=q-1,q}[F_h(S)-F_h(K_c)]$.
Prefixes are the original PRD17 values
of $G_t,G_s$. Coincident flags retain their positions and signs.
The prefix identity follows from the same fixed determinant
identities because its evaluation coefficients sum to zero.
Adding the two equalities proves exactly the NFR7 receiver
\[
 \mathcal M^p
 =\mathcal L_e^p-V_t^p+V_s^p+V_A^p-V_Q^p
 =\Pi_I^p+\tau_I^p-V_Q^p.                              \tag{PRD26}
\]
In particular the independent $V_Q^p$ does not enter the
four-target-copy CPQ allocation identity. It enters through the
proved same-value relative operator PRD19--20.

Finally let $\mathcal C^p$ retain NFR10/EFP31's complete original
compatibility return, including its source-complement projections,
and retain the actual endpoint errors of NFR2:
\[
\begin{gathered}
 a_i^p=2r_I\log\max(1,K_-^p(D_i))+r_I\log c_+^p(D_i),\\
 b_i^p=2r_I\log\max(1,K_+^p(D_i))-r_I\log c_-^p(D_i),\\
 E_p^-=-b_0^p-b_1^p-a_2^p-a_3^p,\qquad
 E_p^+=a_0^p+a_1^p+b_2^p+b_3^p.
\end{gathered}                                                    \tag{PRD27}
\]
NFR5's exact equality is
$\mathcal M^p=\mathcal C^p+\mathcal R\log\omega_{\rm inv}^p
                  +\mathcal E^p$ with
$E_p^-\le\mathcal E^p\le E_p^+$.
Substituting all three lines of PRD21, before estimating any
term, proves the sharpened receiving expression
\[
\begin{gathered}
 \mathcal C^p=\Pi_A^p+\tau_A^p-V_Q^p-\mathcal E^p,\\
 \Pi_A^p+\tau_A^p-V_Q^p-E_p^+
 \le\mathcal C^p\le
 \Pi_A^p+\tau_A^p-V_Q^p-E_p^- .
\end{gathered}                                                    \tag{PRD28}
\]
This preserves the four actual compatibility constants with the
two low signs reversed. It also proves the exact relation
between the complete compatibility return and the proper
pullback correction without dropping their finite error.

All existing source-order intervals therefore receive these
expressions by literal substitution. For example, on EFP29's
already proved finite guards, with $\Delta$ meaning order one
minus order zero,
\[
 r_I(\Phi_j^- -\Phi_k^+)
 \le\Delta(\Pi_I+\tau_I-V_Q)
 \le r_I(\Phi_j^+ -\Phi_k^-).                           \tag{PRD29}
\]
The degree-$k$ constants here are evaluated at PRD22, exactly as
EFP28--29 prescribe, and their target constants remain degree $j$.
For EFP31 its same unchanged full compatibility interval applies
to $\Delta(\Pi_A+\tau_A-V_Q-\mathcal E)$ by PRD28.
No guard is enlarged: EFP28's sufficient $k\ge37$ and actual
QOT root tests, and the separate EOR tests wherever used, remain
those of the source theorem. Outside those comparison guards,
the exact finite matrix identities PRD1--28 retain their original
rank-valid domain.
\endgroup

% Bounded original-action receiving calculation; not a standalone compiled book.
% Accepted782 source SHA256 ef92a4a0f1d9de269a78d698128837ec4ef7c6d82d1c3e9e4c58ef487c9adb64.
% Complete immediate providers: CAI1--15,29--35; NCT1--21; IPR1--14;
% HJR32--35; NSR30--33; NFR1--12; NTG18--29; STF14--18; NFI1--11.
\begingroup
\section{The analytic tail inside the complete original action}
\label{sec:TAC}

This continuation preserves the programme's recorded human origin and the
separate recorded creators and dates of its deposited source notes
\cite{human:globalization2026,human:splitzero2026}. Its new use of the
accepted NTG/STF/NFI identities is the complete original-action return
proved below. The exact packet morphisms extend the saved typed-action
calculation; the full source lineage is recorded in the accompanying
private use ledger. Every new finite Hermitian factorization is proved
here, with the classical Schur references cited where used.

All statements here retain the hypothetical original simple quartet,
the fixed actual period and branch, every conductor coefficient, and the
original source masses.  Write the actual packet degree as $a$ until its
value is specified.  In the current NTG/STF/NFI calculation it is
$a=j=k+4$, not $k$.  Set
\[
 q=q_a=(a+1)^2,\quad q'=q'_a=(a-7)^2,\quad
 m=q-q'-v,\quad N_i=(q-1,q,2q-1,2q)_i,\quad D_i=N_i-v,
 \quad s\in\{1,a\},\quad a=4l+1\ge9.
 \tag{TAC1}
\]
The original source centres are $c_0=a/2$ for the whole numerator
packet and $c_1=a/2-4$ for the divided analytic functionals.  The common
Gamma measure and orthonormal polynomial phases are precisely
\[
\begin{gathered}
 dm_s(y)=\frac{M_s2^{s/2}}{4\pi\Gamma(s/2)}
       |\Gamma(s/4+iy/2)|^2\,dy,\qquad M_s=(2\pi)^{s/2},\\
 p_{0,s}=1,\quad p_{1,s}=y,\quad
 p_{n+1,s}=yp_{n,s}-n(n-1+s/2)p_{n-1,s},\\
 h_{n,s}=M_sn!(s/2)_n,\qquad
 \phi_{n,s,c}(S)=p_{n,s}((S-c)/i)/\sqrt{h_{n,s}}.
\end{gathered}\tag{TAC1a}
\]
The classical Meixner--Pollaczek dictionary is
$p_{n,s}(y)=n!P_n^{(s/4)}(y/2;\pi/2)$; the substitution, including
$dy=2\,d(y/2)$, gives the displayed mass and norms, as calculated
in NTG3 and HPS1--5
\cite[eqs.~(18.19.8)--(18.19.9), (18.22.8)]{human:dlmf18}.  The symbol $m$ in TAC1 is the analytic
coefficient dimension; CAI's quartet multiplicity is exactly one here.
This calculation does not extend the NTG result to a multiple quartet.


\subsection{Both stages of the analytic conductor quotient}

Here are the complete maps from the original packet. Put
\[
\begin{gathered}
 \beta_{a,u,w}=a/2+(2u-a)\delta+i(2w-a)\gamma,\qquad
 \chi_a(S)=\prod_{u,w=0}^{a}(S-\beta_{a,u,w}),\\
 \chi_-(S')=\chi_{a-8}(S'),\quad
 E=\mathbb C[S]/(\chi_a),\quad E_-=\mathbb C[S']/(\chi_-),\\
 E_A(z)=\sum_{r,t=0}^8 a_{rt}e^{\beta_{8,r,t}z},\quad
 \mu_n=E_A^{(n)}(0),\quad v=\operatorname{ord}_0E_A\le80,
 \quad \mu_v\ne0,\\
 (\mathcal T_AP)(S')=\sum_{r,t=0}^8a_{rt}P(S'+\beta_{8,r,t}),\quad
 \overline{\mathcal T}_A:E\longrightarrow E_-,\quad
 [P]\longmapsto[\mathcal T_AP],\\
 L_v=\{[P]:\deg P<v\}\subset E,\qquad
 V_A=\ker\overline{\mathcal T}_A.
\end{gathered}\tag{TAC1b}
\]
The complete coefficients of $E_A$ remain as in the original fixed
period; no repeated contribution is removed. Distinctness of the grid
roots follows from $\delta>0$ and $\gamma>0$ by comparing real and
imaginary parts. For every lower root $\beta_{a-8,u,w}$,
$\beta_{a-8,u,w}+\beta_{8,r,t}=\beta_{a,u+r,w+t}$.
Consequently $\chi_-\mid\mathcal T_A(\chi_a P)$, proving that the
packet arrow in TAC1b is well defined. Its matrix in upper and lower
root-value coordinates is
\[
 C_{(u,w),(i,j)}=
 \begin{cases}a_{i-u,j-w},&0\le i-u,j-w\le8,\\0,&\text{otherwise}.
 \end{cases}\tag{TAC1c}
\]
Indeed evaluation of the complete shift sum gives this entry, with all
contributions retained. For a lower coefficient column $\eta$,
\[
 \mathcal N_{C^{\mathsf T}\eta}(z)
     =E_A(z)\sum_{u,w=0}^{a-8}\eta_{uw}e^{\beta_{a-8,u,w}z}.
\]
If this vanishes, the lower sum vanishes on an open disk on which
$E_A\ne0$, hence identically; its distinct-root Vandermonde then
kills $\eta$. Thus $C^{\mathsf T}$ is injective and $C$ is onto.
In particular $\dim V_A=q-q'$. The polynomial identity
\[
 \mathcal T_AS^n=\sum_{d=0}^n\binom nd\mu_dS'^{n-d}
\]
shows that the exact polynomial kernel is $\mathbb C[S]_{<v}$ and
that the degree of every nonzero image of degree $n\ge v$ is $n-v$,
with leading coefficient multiplied by $\binom nv\mu_v$.
Triangular elimination proves surjectivity onto every retained target
polynomial space. Since $v<q$, $L_v$ injects into $E$, and $L_v\subset V_A$.

For any $N\ge q-1$, put $D=N-v$ and $b=N-q\ge-1$; negative
polynomial degree means the zero space. Define the full relation map
and both quotient spaces by
\[
\begin{gathered}
 \mathcal U_b h=\mathcal T_A(\chi_a h)/\chi_-,\qquad
 X_N=\mathbb C[S']_{\le D}/
                \mathcal T_A(\chi_a\mathbb C[S]_{\le b}),\\
 Y_N=\mathbb C[S']_{\le D-q'}/
                \mathcal U_b\mathbb C[S]_{\le b},\\
 \kappa_N:E/L_v\longrightarrow X_N,\quad
    [P]+L_v\longmapsto[\mathcal T_AP],\\
 \theta_N:V_A/L_v\longrightarrow Y_N,\quad
    [P]+L_v\longmapsto[\mathcal T_AP/\chi_-].
\end{gathered}\tag{TAC1d}
\]
Choose a polynomial representative of degree at most $N$ in these
formulas. Changing it by $\chi_a h$ changes the two images by their
respective displayed relation; changing it by a low polynomial changes
neither image. Surjectivity of $\mathcal T_A$ proves surjectivity of
$\kappa_N$. If its value is zero, subtract the corresponding
$\chi_a h$; the result has degree below $v$, proving injectivity.
The same argument after multiplying a prescribed $Y_N$ representative
by $\chi_-$ proves both surjectivity and injectivity of $\theta_N$.
Every required preimage has degree at most $N$, by the proved exact
degree drop. Thus these are proved isomorphisms, with inverse given by
any such preimage followed by its remainder class modulo $L_v$.

Multiplication by $\chi_-$ and reduction modulo $\chi_-$ give the
commutative diagram of exact sequences
\[
\begin{array}{ccccccccc}
0&\longrightarrow&V_A/L_v&\longrightarrow&E/L_v
 &\xrightarrow{\ \overline{\mathcal T}_A\ }&E_-&\longrightarrow&0\\
 &&\theta_N\downarrow&&\kappa_N\downarrow&&\Vert\\
0&\longrightarrow&Y_N&\xrightarrow{\ \chi_-\ }&X_N
 &\xrightarrow{\ \bmod\chi_-\ }&E_-&\longrightarrow&0 .
\end{array}\tag{TAC1e}
\]
The bottom left map is injective: $\chi_-g=\mathcal T_A(\chi_a h)$
is equivalent to $g=\mathcal U_bh$ in the polynomial integral domain.
Its image is the kernel of the bottom right map, since every polynomial
vanishing at all lower roots is divisible by $\chi_-$. That right map
is onto because $D\ge q-1-v\ge q'-1$. This proves each exactness
claim and the commutativity on actual representatives.
Finally $\mathcal U_b$ is injective: for nonzero $h$ of degree $t$,
its image has degree $q+t-v-q'=m+t$ and leading coefficient
$\operatorname{lc}(h)\binom{q+t}{v}\mu_v\ne0$.
Subtracting the actual polynomial dimensions gives
\[
 \dim(E/L_v)=\dim X_N=q-v,
 \qquad \dim(V_A/L_v)=\dim Y_N=q-v-q'=m.
 \tag{TAC1f}
\]
This displays both stages with their maps and their complete relation
spaces. In the target application $a=j=k+4$, the NTG coefficient
metric has the latter dimension $m=16j-48-v$.

The algebraic pairing identifying its analytic dual is equally exact.
Write $\mathbb V[n,\beta]=\beta^n$ and let $p$ be the monomial
remainder coefficient column. The isomorphism
\[
 \Phi:\mathbb C^q\longrightarrow E^\vee,\qquad
 \Phi(\nu)([P])=\sum_\beta\nu_\beta P(\beta)
                       =(\mathbb V\nu)^{\mathsf T}p
 \tag{TAC1g}
\]
is complex linear; the invertible Vandermonde gives its inverse.
The first $v$ moment conditions say exactly
$\Phi(\mathscr N_v)=\operatorname{ann}(L_v)$.
For $\eta\in\mathbb C^{q'}$,
$\Phi(C^{\mathsf T}\eta)([P])=\sum_{\beta'}\eta_{\beta'}
(\mathcal T_AP)(\beta')$, the actual dual of the onto conductor.
Its image is therefore $\operatorname{ann}(V_A)$: one inclusion
follows by evaluation, and both dimensions are $q'$.
Restriction of functionals induces
\[
 \Psi:\mathscr N_v/\operatorname{im}C^{\mathsf T}
       \xrightarrow{\ \sim\ }(V_A/L_v)^\vee,\qquad
 [\nu]\longmapsto\bigl([P]+L_v\longmapsto\sum_\beta\nu_\beta P(\beta)\bigr).
 \tag{TAC1h}
\]
Adding a low polynomial or a conductor coefficient changes that value
by zero. Its kernel is exactly the stated conductor image. Any
functional on $V_A/L_v$ extends linearly from $V_A$ to $E$ while
still vanishing on $L_v$; its unique $\Phi$ representative lies in
$\mathscr N_v$. This proves surjectivity without a metric choice.
For $F=\mathcal N_\nu/E_A$, analyticity follows from the first $v$
vanishing moments. The exact finite Leibniz identity is
\[
 \Lambda_F(\mathcal T_AP)=\Lambda_{E_AF}(P)
      =\Phi(\nu)([P]),\qquad
 \ell_F(g)=\Lambda_F(\chi_-g),\qquad
 \ell_F(\theta_N([P]+L_v))=\Psi([\nu])([P]+L_v).
 \tag{TAC1i}
\]
For proof on $P=S^n$, expand the complete derivative product as
$\sum_{d=v}^n\binom nd\mu_dF^{(n-d)}(0)$; finite linearity proves
the identity for every retained polynomial. On a relation
$g=\mathcal U_bh$ its value is $\Phi(\nu)([\chi_ah])=0$.
Thus it is a well-defined dual pairing on $Y_N$. A lower exponential
has zero pairing because $\chi_-(\beta')=0$; conversely TAC1h and
the isomorphism $\theta_N$ identify exactly that denominator via
$C^{\mathsf T}$. Every scalar is one in TAC1d--i. The original
$\alpha_A=v!/\mu_v$ remains in the separate primary-coordinate
maps of NFR1 and is not inserted into these unscaled conductor maps.
These arguments prove the CEP3--6,42--43 and NCT3--11 bridge on
its original objects and coordinates.

\subsection{The exact subquotient and its two complementary factors}

Let $\mathbb V$ have entries $\mathbb V[n,\beta]=\beta^n$,
$0\le n<q$, in the original upper-root order.  For
$\nu\in\mathbb C^q$ put $\mathcal N_\nu(z)=\sum_\beta\nu_\beta
e^{\beta z}$.  Define
\[
 \mathscr N_v=\{\nu:(\mathbb V\nu)_n=0\ (0\le n<v)\},
 \quad C=C_a,\quad Z=Z_{\exp,a},\quad
 Z_v=\mathbb V^{-1}P_{<v},\quad
 F=[C^{\mathsf T},Z,Z_v],\quad U=[C^{\mathsf T},Z],
 \tag{TAC2}
\]
where $P_{<v}$ inserts the first $v$ standard coordinate vectors.
NCT6--7 proves that $C^{\mathsf T}$ is injective and that $U$ is a
frame of $\mathscr N_v$, with respective column counts $q'$ and $m$.
The first $v$ moment rows vanish on $U$ and form the identity on $Z_v$.
Thus $F$ is invertible: a relation among its columns first has zero
$Z_v$ coefficient, and then has zero $U$ coefficient by independence.
The spaces and dimensions are therefore exactly
\[
 0\subset\operatorname{im}C^{\mathsf T}\subset\mathscr N_v
       \subset\mathbb C^q,\qquad q=q'+m+v.
 \tag{TAC3}
\]

Let $G_N^{(s)}$ be the full CAI2 remainder-fibre Gram on
$E_a=\mathbb C[S]/(\chi_a)$ with Gamma source order $s$ and centre
$c_0$.  It has the original monomial remainder frame.  Let
\[
 \mathsf E_{N,s}[n,\beta]=\phi_{n,s,c_0}(\beta),\qquad
 K_N^{(s)}=\mathsf E_{N,s}^*\mathsf E_{N,s}.
\]
This is the complete upper observation matrix NCT8, at $0\le n\le N$.
The complex-linear functional associated with $\nu$ sends a remainder
coefficient vector $x$ to $(\mathbb V\nu)^{\mathsf T}x$.
Its squared dual norm is
$(\mathbb V\nu)^*\overline{(G_N^{(s)})^{-1}}(\mathbb V\nu)$:
indeed the representing vector in the Hermitian metric $G_N^{(s)}$ is
$(G_N^{(s)})^{-1}\overline{\mathbb V\nu}$, whose norm has that
value after complex conjugation of its real squared norm.
Computing the same norm in the original orthonormal polynomial basis
gives $\nu^*K_N^{(s)}\nu$.  Equality for every $\nu$ proves
\[
 K_N^{(s)}=\mathbb V^*\overline{(G_N^{(s)})^{-1}}\mathbb V,
 \qquad
 \det K_N^{(s)}=\frac{|\det\mathbb V|^2}{\det G_N^{(s)}}.
 \tag{TAC4}
\]
The ordinary transpose in the functional pairing and the conjugated
inverse in TAC4 are both essential.  This identity is also NCT8a
factored through the original remainder map.  It does not identify
an arithmetic source with a Gamma source.

Suppress $s$ temporarily and define the actual complementary forms
\[
 \begin{split}
 K_{C,N}&=\overline C K_N C^{\mathsf T},\\
 \widehat D_N&=Z^*\{K_N-K_NC^{\mathsf T}K_{C,N}^{-1}
                                      \overline C K_N\}Z,\\
 Q_{v,N}&=Z_v^*\{K_N-K_NU(U^*K_NU)^{-1}U^*K_N\}Z_v.
 \end{split}
 \tag{TAC5}
\]
Here $\widehat D_N$ is exactly NCT8's
$\widehat D_{a,N-v}^{(s)}$, not a new $m$-dimensional quotient.
The metric $K_{C,N}$ is the restriction to every conductor column;
$Q_{v,N}$ is the full attained quotient $\mathbb C^q/\mathscr N_v$
in the fixed values $Z_v$.  In particular its minimization includes
both the conductor columns and all analytic columns.

For clarity, the determinant factorization is proved in these frames.
In $F^*K_NF$, subtract from each of the last $m+v$ column blocks
its $K_N$-orthogonal projection onto $C^{\mathsf T}$, and make the
conjugate row operations.  These triangular operations have determinant
one.  They split off $K_{C,N}$.  Next subtract from the final block
its projection onto the remaining $m$ columns.  This splits off
$\widehat D_N$; minimizing in succession over those two subspaces is
minimizing over their sum $\mathscr N_v$, so the last block is precisely
$Q_{v,N}$.  Consequently
\[
 \det K_{C,N}\det\widehat D_N\det Q_{v,N}
 =|\det F|^2\det K_N
 =\frac{|\det F|^2|\det\mathbb V|^2}{\det G_N^{(s)}}.
 \tag{TAC6}
\]

The fixed determinant factor can also be evaluated in the exact
DNE-selected minor. Let $P_v$ insert the final $q-v$ moment
coordinates, put $J_v=P_v^{\mathsf T}\mathbb V C^{\mathsf T}$,
and retain the first nonzero $q'$-row minor indexed by $\mathcal I$
and the complementary insertion $E_{\mathcal J}$. Then
$Z=\mathbb V^{-1}P_vE_{\mathcal J}$ and
\[
 \mathbb VF=[P_vJ_v,P_vE_{\mathcal J},P_{<v}],\qquad
 |\det F|^2|\det\mathbb V|^2
  =|\det[J_v,E_{\mathcal J}]|^2=|\det(J_v)_{\mathcal I}|^2.
 \tag{TAC6a}
\]
Permute the first $v$ rows to meet the last $v$ columns, then
expand on those identity columns. Expanding the remaining determinant
on the complementary coordinate insertion leaves exactly the stated
minor with a permutation sign. Its absolute square removes the sign
and retains the actual nonzero minor; it is not assigned the value one.

The two outer factors have explicit primal meanings. With
$T=C\mathbb V^{\mathsf T}$, the onto conductor in monomial-to-lower-root
coordinates, and $G=G_N^{(s)}$, put
\[
 G_{C,N}=(TG^{-1}T^*)^{-1},\qquad
 G_{L,N}=P_{<v}^*GP_{<v}.
\]
The original value minimum for $T$ is attained by
$G^{-1}T^*(TG^{-1}T^*)^{-1}$: its composition with $T$ is the
identity and it is $G$-orthogonal to $\ker T$. This proves the
first formula as the full quotient metric. Direct substitution in TAC4
then proves
\[
 K_{C,N}=\overline{G_{C,N}^{-1}},\qquad
 Q_{v,N}=\overline{G_{L,N}^{-1}}.
 \tag{TAC6b}
\]
For the second equality the value map on the dual coefficient packet is
$T_v=P_{<v}^{\mathsf T}\mathbb V$. Its kernel is $\mathscr N_v$
and $T_vZ_v=I_v$, so the attained quotient is
$(T_vK_N^{-1}T_v^*)^{-1}$. Substituting TAC4 gives precisely the
second formula. This proof includes every direction of the full
vanishing-moment kernel.

The middle factor is the dual metric on the actual primal subquotient
$V_A/L_v$, with the restriction of $G$ to $V_A$ followed by its
attained quotient by $L_v$. To prove this, a functional on a Hilbert
subspace extends with minimum norm by composition with the orthogonal
projection. Every other extension differs by a functional on the
orthogonal complement, and their squared norms add. If the original
functional vanishes on $L_v$, that minimum extension still vanishes
on $L_v$; all other admissible extensions differ by
$\operatorname{ann}(V_A)$. This is exactly TAC5's middle minimum
under TAC1g--i. Let $H_{W,N}$ be the primal subquotient Gram transported
by $\theta_N$ to the fixed low monomial frame of $Y_N$, and retain
DNE4's complete pairing matrix $A_\ell$. The complex-linear dual
norm calculation used in TAC4 gives
\[
 \widehat D_N=A_\ell^*\overline{H_{W,N}^{-1}}A_\ell,
 \qquad H_{W,N}=\overline{A_\ell\widehat D_N^{-1}A_\ell^*}.
 \tag{TAC6c}
\]
The pairing matrix is invertible by TAC1h--i. This proves both displayed
formulas and their exact pairing coefficient. NCT22 gives this same
transported Gram. Its native lower-source counterpart replaces
$\widehat D_N$ by $D_{a,N-v}^{(s)}$; TAC9 retains their entire
metric difference.

All positive inverses exist because $K_N>0$ and the stated frames are
independent.  If $v=0$, $Z_v$ and $Q_{v,N}$ are the unique empty maps
and $\det Q_{v,N}=1$.  No limiting inverse is required.
This is the usual Schur factorization, with its complete proof just
given; compare Boyd and Vandenberghe, Appendix A.5.5, p.~650, and
Appendix C.4.1, p.~673 \cite{human:boyd-vandenberghe}.

The native-to-numerator map in this factorization is explicit.
For $D=N-v$, set
\[
 \mathscr Y_D=\{\Lambda_{\mathcal N_\nu/E_A}|_{\mathbb C[S']_{\le D}}
                           :\nu\in\mathscr N_v\},\qquad
 \mathscr K_D=\operatorname{span}\{\operatorname{ev}_{\beta_{a-8,u,w}}\}
                    \subset(\mathbb C[S']_{\le D})^\vee.
\]
Here $\mathscr K_D\subset\mathscr Y_D$ by TAC1c. The complete
finite-jet product identity TAC1i implies injectivity of the map from
$\mathscr N_v$ to $\mathscr Y_D$: zero derivatives through degree
$D$ make every numerator derivative through $N$ zero, and the first
$q$ such derivatives form the invertible upper Vandermonde.
Thus these spaces have dimensions $q-v$ and $q'$ respectively.
For $F_b=\mathcal N_{Ze_b}/E_A$, NCT3--11 and TAC1i give
\[
 [F]\longmapsto[E_AF]:
 \mathscr Y_D/\mathscr K_D
 \xrightarrow{\ \sim\ }
 \mathscr N_v/\operatorname{im}C^{\mathsf T}.
 \tag{TAC7}
\]
On a monomial it is the exact Leibniz identity
\[
 \Lambda_F(\mathcal T_AS^n)
 =\sum_{d=v}^n\binom nd\mu_dF^{(n-d)}(0)
 =(E_AF)^{(n)}(0),\qquad
 (\mathcal T_AP)(S')=\sum_{u,w}a_{uw}P(S'+\beta_{8,u,w}).
\]
The nonzero coefficient $\binom{v+r}{v}\mu_v$ recursively recovers
each derivative of $F$, so this is an isomorphism on the retained finite
jets.  It takes the entire lower-evaluation space to the entire
conductor-multiple space.  Its orthonormal finite-jet matrix is the
ordinary transpose $A_{D,s}^{\mathsf T}$ of NCT3's unscaled conductor
map.  The two quotient metrics on the same $x$ values are the original
$D_{a,D}^{(s)}$ and $\widehat D_{a,D}^{(s)}$; they are not declared equal.

These are exact paired vector-space maps, which need no action
invariance of either intermediate subspace.  The full dual multiplication
action in root coefficients is $\operatorname{diag}(\beta)$ and in
the complete frame $F$ it is
$F^{-1}\operatorname{diag}(\beta)F$, with all its blocks retained.
The conductor's exact action defect is
\[
 \mathcal T_A(SP)-S'\mathcal T_AP
       =\sum_{u,w}a_{uw}\beta_{8,u,w}P(S'+\beta_{8,u,w}),
 \qquad (E_AF)'=E_AF'+E_A'F.
 \tag{TAC8}
\]
Both identities follow term by term from the original shift sum and
the product rule.  The action calculation below uses the whole packet
CAI action; it does not discard TAC8 or manufacture an induced action
on a noninvariant intermediate quotient.


The complete primal action can be recovered from these dual frames,
including every block crossing this flag. This provides an exact action
calculation even when the intermediate subspaces have an action defect.
For any positive original packet Gram $G$ in its monomial frame set
$K=\mathbb V^*\overline{G^{-1}}\mathbb V$, and use all the minima
TAC5 in that $K$. Define
\[
\begin{gathered}
 H=C^{\mathsf T},\qquad
 J=Z-H(H^*KH)^{-1}H^*KZ,\\
 R_\perp=Z_v-U(U^*KU)^{-1}U^*KZ_v,\qquad
 \mathsf B=[H,J,R_\perp],\qquad
 \mathsf D=\operatorname{diag}(K_C,\widehat D,Q_v).
\end{gathered}\tag{TAC8a}
\]
All three block columns are pairwise $K$-orthogonal, as checked in the
proof of TAC6. Thus $\mathsf B^*K\mathsf B=\mathsf D$ and
$\mathsf B^{-1}=\mathsf D^{-1}\mathsf B^*K$. Let
$D_\beta=\operatorname{diag}_{\beta}(\beta)$, and index these block
columns by $u=1,2,3$, of widths $q',m,v$. The full action and every
one of its nine blocks are
\[
 \mathsf A=\mathsf B^{-1}D_\beta\mathsf B,
 \qquad
 \mathsf A_{uv}=\mathsf D_u^{-1}\mathsf B_u^*KD_\beta\mathsf B_v
       \quad(1\le u,v\le3).
 \tag{TAC8b}
\]
This equality follows from the stated inverse, so no off-diagonal
block is suppressed. For example the action of a middle-block vector
$z$ is exactly $(\mathsf A_{12}z,\mathsf A_{22}z,\mathsf A_{32}z)$
in this frame, which explicitly records both components outside the
middle block.

Let $A=M_S$ on $E$. Put $\mathsf P=(\mathbb V\mathsf B)^{-\mathsf T}$, a complete
primal monomial coefficient frame. The pairing is precisely
$\mathsf B^{\mathsf T}\mathbb V^{\mathsf T}\mathsf P=I_q$.
Applying the root action to the functional argument shows that the
primal multiplication matrix in $\mathsf P$ is $\mathsf A^{\mathsf T}$:
for arbitrary functional and primal columns $x,y$, the pairing after
primal multiplication equals $x^{\mathsf T}\mathsf A^{\mathsf T}y$,
which is also the pairing after applying the dual action to $x$.
Moreover direct inversion gives
\[
 \mathsf P^*G\mathsf P=\overline{\mathsf D^{-1}},\qquad
 \mathsf P^*(A^*G+GA-aG)\mathsf P
   =\overline{\mathsf A}\,\overline{\mathsf D^{-1}}
      +\overline{\mathsf D^{-1}}\mathsf A^{\mathsf T}
      -a\,\overline{\mathsf D^{-1}}.
 \tag{TAC8c}
\]
For the first equality, the conjugate inverse of its left side is
$(\mathbb V\mathsf B)^*\overline{G^{-1}}(\mathbb V\mathsf B)
=\mathsf D$. For the second, use $A\mathsf P=\mathsf P\mathsf A^{\mathsf T}$
and substitute the first equality. These computations prove exact
pairing and action compatibility, with all phases and blocks.
Empty blocks are the unique empty maps. Taking $G=G_N^{\rm ar}$
in TAC8a--c gives the complete original CAI arithmetic control;
taking $G=G_N^{(s)}$ gives the original Gamma control. The section
matrices and factors are recomputed in the chosen $G$; the arithmetic
and Gamma forms have not been identified. The explicit arithmetic
volume correction in TAC13 remains in the action below.

\subsection{Whole-volume return with the native factor inserted}

Use the original dual sign operator
$\mathcal R f=-f(q-1)-f(q)+f(2q-1)+f(2q)$ and put
\[
 \begin{gathered}
 \mathcal C_{a,s}=\mathcal R\log\det K_{C,N}^{(s)},\qquad
 \mathcal V_{a,s}=\mathcal R\log\det Q_{v,N}^{(s)},\qquad
 \mathcal T_{a,s}=\mathcal R\log\det D_{a,N-v}^{(s)},\\
 \epsilon_{a,N,s}=\log\det\bigl[(D_{a,N-v}^{(s)})^{-1/2}
          \widehat D_{a,N-v}^{(s)}(D_{a,N-v}^{(s)})^{-1/2}\bigr],
 \qquad \Xi_{a,s}=\mathcal R\epsilon_{a,N,s}.
 \end{gathered}
 \tag{TAC9}
\]
TAC6b identifies the last return with the earlier literal low-polynomial
return, including its orientation:
\[
\begin{aligned}
 \mathcal V_{a,s}
  &=-\mathcal R\log\det G_{L,N}^{(s)}\\
  &=\log\frac{\det G_{L,q-1}^{(s)}\det G_{L,q}^{(s)}}
                 {\det G_{L,2q-1}^{(s)}\det G_{L,2q}^{(s)}}
    =F_{L_a}^{(s)},\qquad L_a=L_v=\mathbb C[S]_{<v}\subset E .
\end{aligned}\tag{TAC9a}
\]
Indeed $G_{L,N}^{(s)}=P_{<v}^*G_N^{(s)}P_{<v}>0$ has positive
real determinant, so TAC6b gives
$\log\det Q_{v,N}^{(s)}=-\log\det G_{L,N}^{(s)}$ at every cutoff.
Applying the displayed four signs proves all equalities. At $v=0$
both logarithms are zero by the empty-determinant convention. Thus the
low factor in this complete packet decomposition is exactly the
original FEX/NSR32 low return used again in TAC17.

Applying $\mathcal R$ to TAC6 cancels the two fixed frame constants
and, because the original CAI volume uses the opposite signs on
$\log\det G_N^{(s)}$, proves the exact whole-volume decomposition
\[
 \boxed{\begin{aligned}
 \mathcal B_a^{(s)}
   &=\mathcal C_{a,s}+\mathcal T_{a,s}+\mathcal V_{a,s}+\Xi_{a,s}\\
   &=\mathcal C_{a,s}+\mathcal T_{a,s}+F_{L_a}^{(s)}+\Xi_{a,s}.
\end{aligned}}
 \tag{TAC10}
\]
The $q'$-dimensional conductor restriction and $v$-dimensional final
quotient are thus the actual complementary factors of the $m$-dimensional
analytic quotient.  They are not CPQ's four complementary returns;
those four decompose a different, four-copy target allocation.

At each original cutoff $K_{N+1}-K_N\succeq0$.  Restriction preserves
this inequality, and taking the minimum of the two quadratic forms on
each identical affine fibre preserves it as well.  Thus $K_{C,N}$,
$\widehat D_N$ and $Q_{v,N}$ increase as positive forms.  Their
$\mathcal R$ returns are nonnegative, by the original telescoping
weights.  This observation does not evaluate their growing-degree
values or change the signed defect $\Xi_{a,s}$.

The finite NCT20 constants retain the exact source, rank, cutoff and
radius. Restore the radius argument suppressed in NCT19, and let
\[
\begin{gathered}
 \mathcal A_A=\{\varrho\in(0,1):g(t)\ne0\text{ for every }|t|=\varrho\},\\
 g(t)=t^{-v}e^{-4w(t)}E_A(w(t)),\qquad w(t)=-i\arctan t,\qquad
 g(0)=(-i)^v\mu_v/v!,\\
 L_i=\inf_{\varrho\in\mathcal A_A}
          \mathcal L_{a,D_i,s,m}^{[1]}(\varrho),\qquad
 U_i=\inf_{\varrho\in\mathcal A_A}
          \mathcal U_{a,D_i,s,m}^{[1]}(\varrho).
\end{gathered}\tag{TAC10a}
\]
Every retained principal part and its full pole order remain in the
NCT16--19 constants. The symbol is analytic in the open unit disk and
$g(0)\ne0$. Its zeros are isolated; on each smaller closed disk they
are finite, so their radii form a countable set. Thus $\mathcal A_A$
is nonempty and contains radii arbitrarily close to one. The displayed
infima are finite and nonnegative, since each fixed admissible radius
has finite nonnegative constants. Attainment of these infima is not
required or asserted. For every admissible radius NCT20 gives
$-\mathcal L(\varrho)\le\epsilon_{a,N_i,s}\le\mathcal U(\varrho)$.
Taking the supremum of the left lower bounds and the infimum of the
right upper bounds proves $-L_i\le\epsilon_{a,N_i,s}\le U_i$.
Define the directed whole-return allowances
\[
 \Xi^-_{a,s}=-U_0-U_1-L_2-L_3,
 \qquad \Xi^+_{a,s}=L_0+L_1+U_2+U_3.
 \tag{TAC11}
\]
Negating the two low inequalities and adding the two high inequalities
proves $\Xi^-_{a,s}\le\Xi_{a,s}\le\Xi^+_{a,s}$. This bounds the
same actual four corrections; it does not assign independent values to
their correlated matrices.

Here is the order of limits for these allowances. For a fixed admissible
$\varrho$, the NCT19 term $4mL_{s,D_i,v}$ is
$O(a^2\log(a+1))$, because $m\le16a$, $s\le a$ and
$D_i\le3a^2$ for $a\ge9$. The exceptional inverse contribution has
rank at most the fixed $H_\varrho$ and is $O_\varrho(a^2)$.
The remaining logarithms obey NCT21's bound
$D_i\log(1/\varrho)+O_\varrho(\log(a+1))$ and have coefficient
at most $2m$. Consequently, for either original source order and
for each endpoint,
\[
 0\le\limsup_{a\to\infty}\frac{L_i}{a^3}
       \le96\log(1/\varrho),\qquad
 0\le\limsup_{a\to\infty}\frac{U_i}{a^3}
       \le96\log(1/\varrho).
 \tag{TAC11a}
\]
For proof, the infimum defining $L_i$ or $U_i$ is bounded above by its
value at this one fixed radius, and the preceding estimates apply with
that radius held fixed throughout the degree limit. Only after taking
this limsup choose admissible radii approaching one. The right sides
tend to zero; therefore $L_i,U_i=o(a^3)=o(aq_a)$ and
$\Xi^-_{a,s},\Xi^+_{a,s}=o(aq_a)$. This last statement concerns the
radius-infimum allowances TAC10a--11. A particular fixed-radius
allowance is retained with its exact finite value whenever it is used;
no zero leading coefficient for that fixed-radius bound is asserted.
No $q'$ rank is inserted into the comparison; its complete metric
remains separately in TAC10.

Now take $a=j=k+4$ and an original $0\le R<q$.  NTG18--25 and
NFI5--6 give the following exact scalar computed by the two tail matrices:
\[
 \begin{split}
 \mathscr T_{a,s,R}
 &:=\sum_{t\in\{q-1,q\}}
              \log\frac{\det\Omega_t}{\det\Omega_R}\\
 &=\frac14\sum_{t\in\{q-1,q\}}\log\det(I_{4m}+M_t)
   =\mathcal T_{a,s}^{>R},\\
 \mathcal T_{a,s}&=\mathcal T_{a,s}^{\le R}+\mathscr T_{a,s,R},
 \qquad 0\le\mathcal T_{a,s}^{\le R}\le\mathfrak E_{s,R}.
 \end{split}
 \tag{TAC12}
\]
The coefficient $1/4$ is forced by NFI6's four copies; each
$\log\det\Omega$ ratio already represents the single native target.
The two endpoint expressions are exactly
$f(2q-1)+f(2q)-2f(q+R)$, so no extra last-row weight is inserted.
At $R=q-1$ the first term is zero, consistently with the empty stack.
IPR2--3 and IPR6 prove the last two statements.  At
$R=\lfloor j^{3/2}/\log j\rfloor$ on its stated range, IPR9 gives
$\mathfrak E_{s,R}=O_{\rm actual}(j^{5/2}+j^3/\log j)=o(kq_k)$.
No value for $\mathscr T_{a,s,R}$ is assigned by this identity.

\subsection{Exact complete action and finite receiving interval}

For this same packet degree $a$, keep every CAI13 argument and define
the actual arithmetic source correction
\[
 \delta_{a,s}^{\rm ar}:=\mathcal B_a^{\rm ar}-\mathcal B_a^{(s)}
 =\begin{cases}
     \delta_a^\sigma,&s=1,\\
     \mathcal H_a+\delta_a^\sigma,&s=a.
   \end{cases}
 \tag{TAC13}
\]
Both equalities follow from CAI5, with no change of measure.
In the literal notation $d_n=V_n^{\rm ar}/V_{n-1}^{\rm ar}$,
CAI12--15 has
\[
 \begin{gathered}
 W_a=\log\frac{\omega_{2q-1,a}^{\rm ar}\omega_{2q,a}^{\rm ar}}
                      {\omega_{q-1,a}^{\rm ar}\omega_{q,a}^{\rm ar}},\\
 P_{a,-}=-\log(1-d_{2q-1})-2\sum_{n=q}^{2q-2}\log(1-d_n),\\
 P_{a,+}=-\log(1-d_q)-\log(1-d_{2q})
                          -2\sum_{n=q+1}^{2q-1}\log(1-d_n),\\
 J_a=2\sum_{N=q-1}^{2q-1}w_N\log\epsilon_N
      =\mathcal B_a^{\rm ar}+W_a-P_{a,-}-P_{a,+},\qquad
 w_{q-1}=w_{2q-1}=1,\quad w_N=2\text{ inside}.
 \end{gathered}
 \tag{TAC14}
\]
The weight sum is $2q$.  The factor $2$ before the weighted logarithms
means the exterior benchmark is $4q\log(\delta q(a+1)/2)$.
Substituting TAC10--13 into this already proved identity gives
\[
 \boxed{\begin{split}
 J_a={}&\mathcal C_{a,s}+F_{L_a}^{(s)}
       +\delta_{a,s}^{\rm ar}+\Xi_{a,s}
       +\mathcal T_{a,s}^{\le R}+\mathscr T_{a,s,R}\\
       &+W_a-P_{a,-}-P_{a,+}.
 \end{split}}
 \tag{TAC15}
\]
Every term has the original $a$-packet and window.  For the current
new matrices $a=j$; TAC15 does not substitute the target matrices at
$j$ for the independent degree-$k$ proper-source metric of NFR1.

The original $W/P$ results are consumed without rederivation.
CAI29 and CAI28 give
$-E_{\rm hi}\le W_a-\mathcal W_q^\sigma\le E_{\rm lo}$ and
$0\le P_{a,-}+P_{a,+}\le\Pi_{h,a}$, with the complete constants of
those formulas, evaluated at $a$ and its $q_a$.
Set the actual retained centre
\[
 A_{a,s,R}=\mathcal C_{a,s}+F_{L_a}^{(s)}
              +\delta_{a,s}^{\rm ar}+\mathscr T_{a,s,R}
              +\mathcal W_q^\sigma.
\]
Adding the four directed bounds, without choosing independent values
for any actual correlated determinant, gives the finite enclosure
\[
 \boxed{\begin{split}
 A_{a,s,R}+\Xi^-_{a,s}-E_{\rm hi}-\Pi_{h,a}
       &\le J_a,\\
 J_a&\le A_{a,s,R}+\Xi^+_{a,s}
                         +\mathfrak E_{s,R}+E_{\rm lo}.
 \end{split}}
 \tag{TAC16}
\]
On ECT's stronger finite guards its already proved penalty cap may
replace $\Pi_{h,a}$; no new penalty theorem is asserted here.
The exact classical-source window remains CAI31's
\[
 \mathcal W_q^\sigma
 =2\log\frac{(4q)!}{(2q)!}-4q\log2+
                         \log\frac{2q-1}{2(4q-1)}.
\]
Equivalently, CAI32--35 supplies $C(q)=4q\log q+(8\log2-4)q$
and its full finite errors $-\log2$, $\eta_q$, $E_{\rm lo}$,
$E_{\rm hi}$ and $\Pi_{h,a}$ in TAC15--16.  The new statement
has not replaced $\delta_{a,s}^{\rm ar}$ or either complementary
volume by an unproved estimate.

\subsection{Why the allocation enters twice and cancels once}

Use the same $a=j$ and order pair that defines the CPQ/STF terms.  Put
\[
 Z_e^X=\mathcal L_e^X+V_b^X+V_{\rm rem}^X+V_s^X+V_A^X,
 \qquad X\in\{\le R,>R\}.
\]
IPR7 and STF18 give $Z_e^X=4S_{K,a}^{X,(s)}$.  The exact
FEX/RQB33/NSR32 correction is
\[
 \begin{gathered}
 \delta_{K,a,s}=\sum_{i=0}^3\sigma_i e_{N_i,s}
                  -F_{L_a}^{(s)}+F_{(K_a+L_a)/K_a}^{(s)},
 \quad \sigma=(1,1,-1,-1),\quad L_a=\mathbb C[S]_{<v},\\
 F_{K,a}^{(s)}=\tfrac14(Z_e^{\le R}+Z_e^{>R})-\delta_{K,a,s}.
 \end{gathered}
 \tag{TAC17}
\]
Here $e_{N_i,s}$ is the actual conductor source Jacobian and both
displayed low metrics remain in their original primal frames.
At source $s=a$, substitution into CAI13/15 gives, term for term,
\[
 \begin{split}
 \Psi_a&=\tfrac14(Z_e^{\le R}+Z_e^{>R})-\delta_{K,a,a}
                  +\mathcal H_a+\delta_a^\sigma,\\
 C_a^{\rm budget}&=\mathcal B_a^{(a)}
       -\tfrac14(Z_e^{\le R}+Z_e^{>R})+\delta_{K,a,a}
                       +W_a-P_{a,-}-P_{a,+},\\
 \Psi_a+C_a^{\rm budget}
     &=\mathcal B_a^{(a)}+\mathcal H_a+\delta_a^\sigma
                       +W_a-P_{a,-}-P_{a,+}=J_a.
 \end{split}
 \tag{TAC18}
\]
Thus the coefficient of every one of the five CPQ tail terms in
$\Psi_a$ is $+1/4$ and in its matching complete budget is $-1/4$.
The two coefficients of $\delta_{K,a,a}$ are $-1$ and $+1$.
They cancel as exact quantities.  TAC15 instead inserts the total
analytic factor with coefficient one into the proved whole-volume
factorization, retaining its two genuinely complementary factors.
These are different exact substitutions, and TAC18 explains their
compatibility without double-counting a boundary term.

For the independent NFR source keep, at its original displaced cutoffs,
\[
 \mathcal M^p=\mathcal L_e^p-V_t^p+V_s^p+V_A^p-V_Q^p,
 \qquad V_A^p-V_Q^p
   =\mathcal R\log\det\bigl[(Q_{\rm inv}^p)^{-1/2}
                A_{\rm inv}^p(Q_{\rm inv}^p)^{-1/2}\bigr].
 \tag{TAC19}
\]
These are NFR7's identities on the same $Z_R$ values.  The degree-$k$
source, original centre, source orders and NFR2 compatibility error
remain.  The present $j$-packet action decomposition does not delete
them or turn a target quotient into that proper-source form.

Finally, CAI11/37 gives on this same complete original arithmetic class
\[
 4q\log L_{h,a}\le J_a,\qquad
 L_{h,a}=\delta q(a+1)/2,\qquad
 L_{h,a}\le\min_N\epsilon_N\le\exp(J_a/(4q)).
 \tag{TAC20}
\]
Combining this with TAC16 yields the explicit necessary inequality
\[
 4q\log L_{h,a}\le A_{a,s,R}+\Xi^+_{a,s}
                            +\mathfrak E_{s,R}+E_{\rm lo}.
\]
The unevaluated terms in this right side are their literal actual
matrices and arithmetic correction.  No smaller numerical upper bound
than the benchmark has been proved here.  The known canonical lower
bounds, and a large weighted mean, do not by themselves contradict
TAC20.  When only a specified complete volume is transported and
$W_a,P_{a,-},P_{a,+}$ keep these arguments, direct subtraction of
TAC14 gives exactly $\widehat J_a-J_a=\widehat{\mathcal B}_a-
\mathcal B_a$.  An analytic-factor substitution qualifies for that
identity only through TAC6--15 with all other factors retained.
\endgroup

\begingroup
\subsection{The same finite joint control in the complete action}
In TAC15--16 put $a=j$ and $R=R_j$ on the already stated
$R_j<q_j$ domain. The original single-copy tail there is exactly
$\mathscr T_{j,s,R}=\tau_{\rm total}$ by FTR9--10 and TAC12.
Consequently the three simultaneous finite estimates retained in
FTR27--28 give, on the same full original arithmetic packet,
\[
\begin{split}
 J_j\le{}&\mathcal C_{j,s}+\mathcal V_{j,s}
     +\delta_{j,s}^{\rm ar}+\mathcal W_{q_j}^{\sigma}
     +\mathfrak U_R+\Xi_{j,s}^{+}
     +\mathfrak E_{s,R}+E_{\rm lo}.                      \tag{FTR31}
\end{split}
\]
To prove this, substitute the exact TAC15 equality, then apply
$\mathscr T_{j,s,R}\le\mathfrak U_R$ from FTR28,
$\Xi_{j,s}\le\Xi_{j,s}^+$ from TAC11,
$0\le\mathcal T_{j,s}^{\le R}\le\mathfrak E_{s,R}$ from TAC12,
$W_j\le\mathcal W_{q_j}^{\sigma}+E_{\rm lo}$ from CAI29,
and $P_{j,-}+P_{j,+}\ge0$ from the actual CAI penalties.
Every term has coefficient one in TAC15. The exact two complementary
packet volumes and arithmetic correction remain on the right of FTR31.
The finer attained value $\tau_{\rm total}$ may always be retained
in place of its upper bound, giving TAC16's sharper exact-data centre.

The four CPQ complements enter the allocation through FTR10 and
TAC17 with coefficient $1/4$; at TAC18's stated order $s=j$,
they enter the matching budget with coefficient $-1/4$.
They are therefore retained with both
coefficients before cancellation; they do not replace either packet
complement in FTR31. The independently computed relative proper-source
returns remain PRD21 and PRD28 at their degree-$k$ cutoffs. This
identifies the actual receiving maps for all three uses of the same
tail calculation, with their original domains, values and coefficients.
\endgroup
